\section{Experimental Setup}
\label{sec:experimental-setup}

\subsection{Production Evaluation Data}

The evaluated system is a production analytics agent designed to answer
natural-language questions about data. It combines language-model reasoning
with tools for retrieving and analyzing information, and is evaluated
repeatedly as its models, prompts, tools, and surrounding system change.

The benchmark consists of human-curated questions about data written by internal analysts; neither the questions nor the recorded evaluation outcomes contain user data or personal information. Each recorded outcome is a binary pass/fail verdict produced by an automated grader. We
retain evaluation runs with valid outcomes for at least 80\% of the benchmark
and compute each run's reference pass rate over its valid outcomes. 
This yields 287 calibration runs from Day 1 through Day 28, and 287 test runs from Day 29 through Day 52 (Figure~\ref{fig:overview}).


\subsection{Temporal Protocol and Method Configurations}

\paragraph{Random sampling.}
We evaluate $k\in\{10,\allowbreak20,\allowbreak30,\allowbreak50,\allowbreak75,\allowbreak100,\allowbreak150,\allowbreak200,\allowbreak250,\allowbreak300,\allowbreak350,\allowbreak400\}$, averaging
results over 20 random seeds.

\paragraph{Historical caching.}
Caching uses a rolling seven-day window of prior routine runs, requires at least
five historical trials per eligible question, and is rebuilt daily for reuse on
that day’s evaluations. We sweep $\epsilon$ from 0.00 to 0.30 in increments of
0.01, caching a question when
Eq.\ref{equation:cache-eligible} is satisfied.

\paragraph{Fixed-subset evaluation.}
We construct each subset once from all calibration runs at the same computational budgets $k$ used for random sampling and apply it unchanged to every test run. IRT models for selection and
p-IRT reconstruction are also fitted on all calibration runs. To estimate the
gp-IRT bias term $\hat{b}^2$, we split calibration runs in half, fit an auxiliary
IRT model on the first half, and measure reconstruction error on the second by
revealing half of each run’s available outcomes. For every multidimensional 2PL
fit, we select $d\in\{2,5,10,15\}$ by log likelihood on a fixed 15\% sample of
observed calibration cells.


\paragraph{Adaptive testing.}
We evaluate adaptive testing at the same $k$. Item parameters are fitted on all calibration runs and held fixed during
question selection for each test run. We compare Rasch adaptive testing with the multidimensional 2PL
variant to assess whether modeling multiple latent abilities improves the
accuracy--execution tradeoff.